\chapter{Critical coefficients at a collision}
\label{cycoll:ch}

An extension of two commuting representations must itself satisfy
the moment equation. When their eigenvalues meet, that equation can
lose rank. Its derived zero fiber then retains classes that disappear
if one first passes to ordinary polynomial functions. We begin by
computing those classes on a two-dimensional flag chart.

A second phenomenon appears when three positions meet. A local
system can extend cleanly across each punctured pair-collision
divisor and still have cohomology at their common intersection.
An explicit cubic cover realizes this situation. Its critical
direct image contains a point constituent that its intermediate
extension does not contain. Both examples explain why a comparison
of Hall coefficients with local fields must specify the coefficient
complex and its extension across collisions.

\section{The moment equation on a flag}
\label{cycoll:sec:moment}

Let $V=\mathbb C^2$, with its first coordinate line.
Consider the matrices and the framing maps
$i:\mathbb C\to V$, $j:V\to\mathbb C$ given by
\[
 X=\begin{pmatrix}x_1&B_0\\0&x_2\end{pmatrix},\qquad
 Y=\begin{pmatrix}y_1&D_0\\0&y_2\end{pmatrix},\qquad
 i=\binom{i_1}{i_2},\qquad j=(1,0).
\]
The equation is $[X,Y]+ij=0$. The framing is \emph{cocyclic} when
no nonzero subspace invariant under both matrices lies in $\ker j$.
On this slice that condition is $(B_0,D_0)\ne(0,0)$: the only
possible such subspace is the second coordinate line.

Direct multiplication gives the four entries of the moment map:
\begin{equation}\label{cycoll:eq:four-moments}
 i_1,\qquad (x_1-x_2)D_0-(y_1-y_2)B_0,\qquad i_2,\qquad 0.
\end{equation}
The last entry is a zero equation. It still contributes a generator
to the specified derived zero fiber.

We use cohomological grading: the differential has degree $1$,
and $H^q(E[s])=H^{q+s}(E)$. For functions $f_1,\ldots,f_r$ on
an affine scheme with ring $A$, their Koszul algebra is
$A\otimes\Lambda(\xi_1,\ldots,\xi_r)$, with
$|\xi_i|=-1$ and $d\xi_i=f_i$. This algebra specifies the derived
zero fiber of the given list of functions. Replacing the list by
equations for the same ordinary zero set can change that fiber.

The pairs for $i_1$ and $i_2$ in
\eqref{cycoll:eq:four-moments} are contractible. Indeed
$A[i]\otimes\Lambda(\xi)$, with $d\xi=i$, retracts onto $A$:
the projection sets $i=\xi=0$, and the degree-$-1$ homotopy is
\[
 H(i^m)=\xi i^{m-1}\quad(m\geq1),\qquad
 H(1)=H(\xi i^m)=0.
\]
Tensoring these contractions with the other factors leaves
\[
 d\eta=(x_1-x_2)D_0-(y_1-y_2)B_0,\qquad d\chi=0,
 \qquad |\eta|=|\chi|=-1.
\]

On $B_0\ne0$, the remaining diagonal change of basis fixes $B_0=1$,
with the corresponding invertible rescaling of the equation generators.
Take the derived restriction to the curve $y_1=y_2=0$ in the plane.
The Koszul complex is termwise free over the eigenvalue base, so
derived restriction is computed by substitution in that complex.
The polynomial model is defined over $\mathbb Q$. After base change
to a nonzero commutative $\mathbb Q$-algebra $R$, it is
\begin{equation}\label{cycoll:eq:moment-complex}
 K=R[x_1,x_2,D_0]\otimes\Lambda(\eta,\chi),\qquad
 d\eta=(x_1-x_2)D_0,\qquad d\chi=0.
\end{equation}
The three polynomial variables have cochain degree zero. We retain
only this grading in
the following calculations.

The variables $x_i,y_i$ are plane eigenvalues. They are distinct
from the Chern roots used in an equivariant Hall formula. An
identification of the moment complex with a Hall direct-image
complex requires the actual support maps, orientation and shift.

\begin{proposition}\label{cycoll:prop:polynomial-fibre}
Put $u=x_1-x_2$ and $w=x_2$. Then
\begin{equation}\label{cycoll:eq:polynomial-cohomology}
 H^*(K)=R[u,w,D_0]/(uD_0)\otimes\Lambda(\chi).
\end{equation}
After inverting $u$, the complex retracts onto
$R[u,u^{-1},w]\otimes\Lambda(\chi)$ with zero differential.
Its derived fiber over $u=0$ is instead
\begin{equation}\label{cycoll:eq:derived-fibre}
 K\otimes^L_{R[u]}R
   =R[w,D_0]\otimes\Lambda(\eta,\chi),\qquad d=0.
\end{equation}
In particular the collision fiber has a nonzero class $\eta\chi$
in degree $-2$.
\end{proposition}
\begin{proof}
Multiplication by $uD_0$ is injective on $R[u,w,D_0]$: it sends
distinct polynomial monomials to distinct monomials without
changing their coefficients. This remains true if $R$ has
nilpotents. Thus the two-term complex with generator $\eta$ has
only its degree-zero cokernel. Tensoring with the closed exterior
factor $\Lambda(\chi)$ proves
\eqref{cycoll:eq:polynomial-cohomology}.

After localization, let $p$ set $D_0=\eta=0$ and let $i$ be the
inclusion. Define
\[
 H(D_0^m\chi^e)=u^{-1}\eta D_0^{m-1}\chi^e
 \quad(m\geq1,\ e=0,1),\qquad
 H(\chi^e)=H(\eta D_0^m\chi^e)=0.
\]
The identity $dH+Hd=1-ip$ follows on each monomial. It retains
the closed generator $\chi$.

Finally $K$ is bounded and termwise free over $R[u]$. It therefore
computes its derived tensor product with $R$. Setting $u=0$ makes
the differential vanish and proves
\eqref{cycoll:eq:derived-fibre}.
\end{proof}

Ordinary restriction of the algebra in
\eqref{cycoll:eq:polynomial-cohomology} would miss $\eta$ and
$\eta\chi$. They are the additional derived-fiber classes. The
closed class $\chi$ is present even before this restriction; it
comes from the zero entry in the original moment map.

\section{What polynomial tests forget}
\label{cycoll:sec:projection}

Complete the polynomial variables in $K$ at $(x_1,x_2,D_0)$.
Its continuous $R$-linear dual, with $R$ discrete, consists of
functionals depending on finitely many Taylor coefficients. Denote
it by $K^\circ$, with differential
\[
 (d^\circ\phi)(k)=(-1)^{|\phi|+1}\phi(dk).
\]
It lies in degrees $0,1,2$. For $m,n,l\geq0$, write
\[
 v_{mnl}(F)=
 \left.\partial_{x_1}^m\partial_{x_2}^n\partial_{D_0}^lF\right|_0.
\]
The notation $\eta^\vee v$, $\chi^\vee v$ or
$(\eta\chi)^\vee v$ means extraction of the indicated exterior
coefficient followed by $v$. These are linear dual coordinates;
no multiplication of dual odd generators is being asserted.

Give $P_2=R[x_1,x_2]^{\mathfrak S_2}$ the topology induced by
the $(x_1,x_2)$-adic topology. Its continuous dual is
identified with the free $R$-module $W_2$ on the commutative
symbols $J_mJ_n=J_nJ_m$, by
\[
 \langle J_mJ_n,F\rangle
   =\left.\partial_{x_1}^m\partial_{x_2}^nF\right|_0.
\]
To see that the pairing is nondegenerate, use the symmetric
monomial basis. The diagonal entries are $m!n!$, including when
$m=n$, and these are units in $R$. Thus restriction along the
inclusion $P_2\hookrightarrow K$ defines a cochain map
\[
 \Theta:K^\circ\longrightarrow W_2[0].
\]

\begin{proposition}\label{cycoll:prop:lost-classes}
The map $\Theta$ is surjective and has the formula
\begin{equation}\label{cycoll:eq:projection}
 \Theta(v_{mnl})=
 \begin{cases}J_mJ_n,&l=0,\\0,&l>0,\end{cases}
 \qquad
 \Theta(\eta^\vee v)=\Theta(\chi^\vee v)
       =\Theta((\eta\chi)^\vee v)=0.
\end{equation}
It kills the nonzero classes
\begin{equation}\label{cycoll:eq:lost}
 [v_{001}]\in H^0(K^\circ),\qquad
 [\chi^\vee v_{000}]\in H^1(K^\circ).
\end{equation}
It intertwines the transpose of $\partial_{x_1}+\partial_{x_2}$
with $T(J_mJ_n)=J_{m+1}J_n+J_mJ_{n+1}$. In particular translation
compatibility does not make $\Theta$ a quasi-isomorphism.
\end{proposition}
\begin{proof}
Restriction to symmetric polynomials gives
\eqref{cycoll:eq:projection}, and the images span $W_2$.
The transpose of multiplication by $(x_1-x_2)D_0$ is
\[
 L v_{mnl}=ml\,v_{m-1,n,l-1}-nl\,v_{m,n-1,l-1};
\]
a term with zero integer coefficient is zero. The dual
differential is consequently
\begin{align*}
 d^\circ v&=-\eta^\vee Lv,&
 d^\circ(\eta^\vee v)&=0,\\
 d^\circ(\chi^\vee v)&=(\eta\chi)^\vee Lv,&
 d^\circ((\eta\chi)^\vee v)&=0.
\end{align*}
Both $Lv_{001}$ and $Lv_{000}$ vanish. The first class in
\eqref{cycoll:eq:lost} has no incoming boundary and evaluates to
$1$ on $D_0$. Every boundary of degree one has only an
$\eta^\vee$ component, so the second class, which evaluates to
$1$ on $\chi$, is not a boundary either.

The simultaneous derivative commutes with $d$, since it kills
$x_1-x_2$. Its transpose sends $v_{mnl}$ to
$v_{m+1,n,l}+v_{m,n+1,l}$. Formula
\eqref{cycoll:eq:projection} now gives the asserted translation
identity, in every cochain degree.
\end{proof}

The first lost class differentiates in the moment coordinate.
The second is dual to a relation direction. The argument persists
under any scalar specialization of $R$: both detecting evaluations
are equal to $1$. A degree-preserving map to a complex concentrated
in degree zero cannot retain the second class. A geometric direct
image may change the complex and shift; restriction of polynomial
tests does neither.

\section{The separated coefficient and its homotopy}
\label{cycoll:sec:separated}

To compare with a spectral coefficient, specialize the rational
moment model to
\[
 R_4=\mathbb Q(a,b)[\epsilon]/(\epsilon^4),\qquad h=a+b,
 \qquad B=R_4[x,y]((s^{-1})),\qquad Z=s+y-x.
\]
The coordinate $Z$ is invertible in $B$. Expand fractions first
in $\epsilon$ and then at $s=\infty$. The scalar
\begin{equation}\label{cycoll:eq:euler}
 E(Z)=\frac{Z(Z+\epsilon h)}{(Z+\epsilon a)(Z+\epsilon b)}
      =1-\frac{ab\epsilon^2}{Z^2}
           +\frac{abh\epsilon^3}{Z^3}
\end{equation}
is exact at this precision. Define
\[
 K_s=B[D_0]\otimes\Lambda(\eta,\chi),\qquad
 d\eta=-ZD_0,\qquad d\chi=0.
\]
Substitution $x_1=x-s$, $x_2=y$ defines a cochain map
$\sigma_s:K\to K_s$, with $K$ now taken over $R_4$.
Multiplying its output by $E(Z)$ gives the coefficient map
\begin{equation}\label{cycoll:eq:lift}
 \mathcal L_s=E(Z)\sigma_s.
\end{equation}
This map is linear. Its value on $1$ is $E(Z)$, so it is not
asserted to be a unital algebra homomorphism.

Let $p:K_s\to B\otimes\Lambda(\chi)$ set $D_0=\eta=0$,
and let $i$ be the inclusion. With the exterior factors ordered
as $\eta\chi$, put
\begin{equation}\label{cycoll:eq:homotopy}
 H(D_0^m\chi^e)=-\frac{\eta D_0^{m-1}\chi^e}{Z}
 \quad(m\geq1,\ e=0,1),\qquad
 H(\chi^e)=H(\eta D_0^m\chi^e)=0.
\end{equation}

\begin{proposition}\label{cycoll:prop:lift}
The map $\mathcal L_s$ preserves the two odd coefficient
directions and commutes with the simultaneous derivative in
the two roots. The homotopy satisfies
\begin{equation}\label{cycoll:eq:contraction}
 dH+Hd=1-ip,\qquad
 d(EH)+(EH)d=E(Z)(1-ip).
\end{equation}
For a symmetric polynomial $F$,
\[
 p\mathcal L_s(F)=E(s+y-x)F(x-s,y).
\]
The first corrections to the homotopy are
\begin{equation}\label{cycoll:eq:weighted-homotopy}
 (EH)(D_0)=-\frac{\eta}{Z}
          +\frac{ab\epsilon^2\eta}{Z^3}
          -\frac{abh\epsilon^3\eta}{Z^4}.
\end{equation}
\end{proposition}
\begin{proof}
Substitution takes $(x_1-x_2)D_0$ to $-ZD_0$, and multiplication
by the closed scalar $E(Z)$ commutes with the differential.
Both maps preserve multiplication by $\chi$. The simultaneous
derivative kills $Z$, so it commutes with $\mathcal L_s$.

For $m\geq1$, $dH$ is the identity on $D_0^m\chi^e$ and
$Hd$ vanishes there. On $\eta D_0^m\chi^e$, the two roles
are reversed. Both vanish on $1,\chi$. These monomials prove
the first identity in \eqref{cycoll:eq:contraction}.
Multiplication by $E$ proves the second. Projection gives the
polynomial formula, and \eqref{cycoll:eq:euler} gives
\eqref{cycoll:eq:weighted-homotopy}.
\end{proof}

Thus the scalar coefficient also determines a correction to the
moment homotopy. Retaining $E$ in the coefficient map while
omitting it from the homotopy would leave the defect
\begin{equation}\label{cycoll:eq:euler-defect}
 \bigl(dH+Hd-E(Z)(1-ip)\bigr)(D_0)
   =\frac{ab\epsilon^2D_0}{Z^2}\pmod{\epsilon^3}.
\end{equation}
At $a=b=Z=D_0=1$, its $\epsilon^2$ coefficient is $1$.

\section{Principal parts on the collision}
\label{cycoll:sec:supported}

Return to the coefficient ring $R$ of
\eqref{cycoll:eq:moment-complex}. The localized calculation
divides by $u=x_1-x_2$. Coefficients
supported at $u=0$ have a different action of $u$. To compute it,
define the principal-parts module
\[
 P_u=R[u,u^{-1}]/R[u]
     =\bigoplus_{p\geq1}R\,u^{-p}.
\]
Multiplication by $u$ kills $u^{-1}$ and takes $u^{-p}$ to
$u^{-(p-1)}$ for $p>1$. Define $P_v$ in the same way, and set
\[
 \mathcal T=R[w]\otimes_R P_u\otimes_R P_v,
 \qquad \mathcal T_1=\ker(u:\mathcal T\to\mathcal T).
\]
Here $v$ is a second normal coordinate and $w$ is a common
position coordinate. These are algebraic direct sums of
principal parts. In particular every element has finite pole
order in each normal coordinate. We make no assertion about
an ambient product or an adic completion of this module.

Tensor the moment complex with $\mathcal T$ over its coefficient
ring in $u,w$. The result is
\begin{equation}\label{cycoll:eq:supported-complex}
 M=\mathcal T[D_0]\otimes\Lambda(\eta,\chi),\qquad
 d\eta=uD_0,\qquad d\chi=0.
\end{equation}
This is a complex of modules; principal parts are not being
given a quotient-ring multiplication.

\begin{proposition}\label{cycoll:prop:supported}
The cohomology of \eqref{cycoll:eq:supported-complex} is
\begin{align}
 H^0(M)&=\mathcal T,\notag\\
 H^{-1}(M)&=\eta\mathcal T_1[D_0]\oplus\chi\mathcal T,
       \label{cycoll:eq:supported-cohomology}\\
 H^{-2}(M)&=\eta\chi\mathcal T_1[D_0].\notag
\end{align}
The copies of $\mathcal T$ in degrees $0,-1$ are represented
at $D_0=0$. In particular
\[
 [\eta u^{-1}v^{-1}],\qquad[\eta\chi u^{-1}v^{-1}]
\]
are nonzero. There is no deformation retraction of $M$ onto
$\mathcal T\otimes\Lambda(\chi)$ with zero differential.
\end{proposition}
\begin{proof}
Multiplication by $u$ is surjective on $\mathcal T$. Therefore
the image of $uD_0$ on $\mathcal T[D_0]$ is
$D_0\mathcal T[D_0]$, and its kernel is $\mathcal T_1[D_0]$.
These give the cokernel and kernel of the two-term $\eta$
complex. Tensoring with the closed factor $\Lambda(\chi)$
proves \eqref{cycoll:eq:supported-cohomology}.

Explicitly, $\eta u^{-p-1}v^{-r}D_0^{m-1}$ is a primitive
of $u^{-p}v^{-r}D_0^m$ for $m\geq1$. By contrast
$\eta u^{-1}v^{-r}D_0^m$ is closed. A boundary in degree $-1$
has only a $\chi$ component, and there are no boundaries in
degree $-2$. This proves the two stated nonvanishing claims.
The proposed retract has no degree-$-2$ cohomology, whereas
$M$ does. A deformation retraction is therefore impossible.
\end{proof}

The coefficient derivative $\partial_u$ alone is not a
connection on this moment complex. Holding $D_0$ fixed gives
\[
 [\partial_u,d]\eta=D_0.
\]
A horizontal comparison must supply the additional action on
the moment factors or the connection after geometric direct
image. The failure of the localized contraction proves a
statement about that contraction and recipient. It leaves open
a comparison retaining the full supported complex.

\section{The link of a triple collision}
\label{cycoll:sec:link}

We now use bounded constructible complexes of finite-dimensional
$k$-vector spaces on complex algebraic varieties with their
analytic topology. The field $k$ has characteristic zero. All
sheaf operations are derived. For an open inclusion $j$, write
$j_!$ for extension by zero and $Rj_*$ for right direct image.
For a closed point inclusion $i$, $i^*$ is its ordinary stalk
and $i^!$ its costalk, or extraordinary restriction. These two
restrictions detect different supported maps.

For three positions put $u=z_1-z_2$, $v=z_2-z_3$, and remove
the common-position coordinate. The normal surface and the
separated open are
\begin{equation}\label{cycoll:eq:three-lines}
 X=\mathbb C^2,\qquad D=\{uv(u+v)=0\},\qquad U=X\setminus D.
\end{equation}
Write $j:U\hookrightarrow X$ and $i:\{0\}\hookrightarrow X$.
The line $u+v=0$ away from $0$ is the collision $z_1=z_3$
with $z_2$ distinct. It is one of the three branches, and its
coefficient must be included before studying their common
intersection.

The coordinates $(u,v)=(t,tr)$ give
\[
 U\simeq\mathbb C^*\times
       \bigl(\mathbb C\setminus\{0,-1\}\bigr).
\]
The second factor retracts onto a bouquet of two circles.
Let $c$ be the positive loop in the first factor, and let $a,b$
be based positive meridians around $0,-1$ in the second. Then
$\pi_1(U)=\mathbb Z\times F(a,b)$, where $F(a,b)$ is the free
group on $a,b$. With compatible base paths, the remaining
branch meridian is conjugate to $c(ab)^{-1}$.

A finite-dimensional representation $M$ of this group has
invertible actions $C,A,B$, with
$CA=AC$ and $CB=BC$. The operators $A,B$ need not commute.
Let $\mathcal L_M$ be the corresponding local system. Define
\begin{equation}\label{cycoll:eq:link-complex}
 K(M)=\left[M\xrightarrow{d^0}M^3\xrightarrow{d^1}M^2\right]
 \quad\text{in degrees }0,1,2,
\end{equation}
where
\begin{align*}
 d^0m&=((C-1)m,(A-1)m,(B-1)m),\\
 d^1(p,q,r)&=((C-1)q-(A-1)p,\ (C-1)r-(B-1)p).
\end{align*}
The commutation with $C$ gives $d^1d^0=0$.

\begin{proposition}\label{cycoll:prop:link}
There are natural identifications
\[
 R\Gamma(U,\mathcal L_M)\simeq K(M)
       \simeq i^*Rj_*\mathcal L_M.
\]
If $C-1$ is invertible, this complex is contractible. The same
formulas apply to bounded complexes of representations after
totalization with the internal differential and its Koszul signs.
\end{proposition}
\begin{proof}
The universal cover of the angular graph is a tree. Its cellular
resolution has boundary $(a-1,b-1)$. The radial circle has
boundary $c-1$. Tensoring these two free resolutions and applying
$\operatorname{Hom}(-,M)$ gives
\eqref{cycoll:eq:link-complex}, including its minus signs.
The product of the circle and graph is a classifying space for
$\mathbb Z\times F(a,b)$, so this complex computes the stated
cohomology. Positive real dilation identifies the complement in
a sufficiently small ball, up to homotopy, with the global
complement. This proves the stalk assertion.

When $C-1$ is invertible, the maps
\[
 h^1(p,q,r)=(C-1)^{-1}p,\qquad
 h^2(y,z)=(0,(C-1)^{-1}y,(C-1)^{-1}z)
\]
satisfy $dh+hd=1$. Commutation with $A,B$ proves the identity
in degree one; the other two degrees follow directly.
\end{proof}

For the constant local system both differentials vanish, giving
Betti numbers $(1,3,2)$. The complement of two independent
coordinate axes instead has Betti numbers $(1,2,1)$. Thus the
third collision branch changes the link even with constant
coefficients.

Noncommuting monodromy already occurs in dimension two. Take
\[
 C=I,\qquad
 A=\begin{pmatrix}1&1\\0&1\end{pmatrix},\qquad
 B=\begin{pmatrix}1&0\\1&1\end{pmatrix}.
\]
The common invariant space is zero and both differentials have
rank two. The cohomology dimensions are $(0,2,2)$. This example
has invariant vectors around individual branches; it will be
useful to compare it with coefficients having no such invariants.

\section{Clean branches and a surviving point}
\label{cycoll:sec:clean}

We say that a local system extends \emph{cleanly} across a
stratum if the canonical map $j_!\mathcal L\to Rj_*\mathcal L$
is an isomorphism near that stratum. At a smooth branch of $D$,
the normal link is a circle. For scalar monodromy $\lambda$
its cochain complex is $[k\xrightarrow{\lambda-1}k]$.
Thus $\lambda\ne1$ gives clean extension along the punctured
branch.

Take a rank-one coefficient with $C=1$, $A=2$ and $B=3$.
Its three branch monodromies are $2,3,1/6$, so each punctured
branch is clean. At their intersection the link is
\[
 d^0m=(0,m,2m),\qquad d^1(p,q,r)=(-p,-2p).
\]
It has $H^0=0$ and one-dimensional groups
$V_1=H^1K$, $V_2=H^2K$. Representatives are $(0,0,1)$
in degree one and $(0,1)$ modulo $k(1,2)$ in degree two.

\begin{proposition}\label{cycoll:prop:clean-defect}
For this coefficient there is a distinguished triangle
\begin{equation}\label{cycoll:eq:clean-triangle}
 j_!\mathcal L\longrightarrow Rj_*\mathcal L
       \longrightarrow i_*K\longrightarrow j_!\mathcal L[1].
\end{equation}
Moreover $i^!j_!\mathcal L\simeq K[-1]$, with
one-dimensional cohomology in degrees $2,3$. The same conclusions
hold for $C=1$, $A=a$, $B=b$ whenever $a,b,ab\ne1$.
\end{proposition}
\begin{proof}
The cone of $j_!\mathcal L\to Rj_*\mathcal L$ vanishes on
$U$ and on every punctured branch. It is supported at $0$,
where its stalk is $K$. This identifies the cone with $i_*K$.
The localization adjunction gives $i^!Rj_*\mathcal L=0$;
applying $i^!$ to the triangle gives the shift $K[-1]$.
For the scalar family in the statement, $d^0$ and $d^1$ both
have rank one. The branch monodromies are $a,b,(ab)^{-1}$,
all different from $1$.
\end{proof}

Consequently clean pairwise collisions need not imply clean
extension at a triple collision. Among rank-one coefficients,
the locus with $H^1K\ne0$ is exactly $C=1$. Away from that
locus the radial contraction applies. On it the angular
graph has nonzero first cohomology. In branch coordinates,
the resonance equation is that the product of the three
monodromies equals $1$.

The same phenomenon allows noncommuting monodromy. For
\begin{equation}\label{cycoll:eq:clean-matrices}
 C=I,\qquad
 A=\begin{pmatrix}2&1\\0&3\end{pmatrix},\qquad
 B=\begin{pmatrix}3&0\\1&2\end{pmatrix},
\end{equation}
the first two branch matrices have eigenvalues $2,3$, while
$(AB)^{-1}$ has eigenvalues $1/9,1/4$. Every branch is clean.
The invertibility of $A-I$ makes both link differentials
have rank two, again giving dimensions $(0,2,2)$.
Here $AB\ne BA$. Thus the point cohomology is not an artifact
of requiring scalar monodromy.

\section{A cubic cover and its intermediate extension}
\label{cycoll:sec:cubic}

There is an algebraic realization of the scalar resonance.
Adjoin a primitive cube root of unity $\zeta$ to $\mathbb Q$
and consider the finite map
\begin{equation}\label{cycoll:eq:cubic}
 p:Z=\{w^3=uv(u+v)\}\longrightarrow X.
\end{equation}
Over $U$ it is an etale cover with deck group $\mu_3$.
A nontrivial character summand of its constant direct image
has $C=1$, $A=B=\zeta$, after choosing the character.
Indeed, with $(u,v)=(t,tr)$, the variable $y=w/t$ satisfies
$y^3=r(r+1)$. A radial loop fixes the character, and each
of the three branch loops has monodromy $\zeta$.

The projective normalization of $y^3=r(r+1)$ has genus one:
its degree-three map to $\mathbb P^1$ is fully ramified at
$0,-1,\infty$, so Riemann--Hurwitz gives
$2g-2=-6+3(3-1)=0$. Over $\mathbb Q$, the two nontrivial
characters together have matrices
\[
 C=I,\qquad A=B=J,\qquad
 J=\begin{pmatrix}0&-1\\1&-1\end{pmatrix}.
\]
Since $J^3=I$ and $J-I$ is invertible, the rational link
has dimensions $(0,2,2)$. Each character has dimensions
$(0,1,1)$ after adjoining $\zeta$.

To describe the extension across $0$, recall the perverse
normalization. For a stratum $S$ of complex dimension $s$,
a constructible complex $P$ is perverse if its restriction
to $S$ has no cohomology above $-s$ and its extraordinary
restriction has none below $-s$. A local system on a smooth
complex $d$-fold becomes perverse after shift by $d$.
The intermediate extension $j_{!*}\mathcal L[d]$ is the
perverse extension with no subobject or quotient supported
on the complement. Equivalently, it is the image of
$j_!\mathcal L[d]\to Rj_*\mathcal L[d]$ in the perverse
heart, with perverse truncations when needed. We use the
shifted convention of \cite[\S15.6]{cycoll:Goresky}.

Retain a rank-one coefficient with $C=1$ and all three
branch monodromies nontrivial. Let $K$ be its link and
$V_i=H^iK$ for $i=1,2$. Define
\begin{equation}\label{cycoll:eq:middle}
 F=\operatorname{fib}\left(Rj_*\mathcal L
             \longrightarrow i_*V_2[-2]\right),
 \qquad \mathcal I=F[2].
\end{equation}
The arrow is point restriction followed by the canonical
top-cohomology quotient $K\to H^2(K)[-2]$.

\begin{theorem}\label{cycoll:thm:middle}
The complex $\mathcal I$ is $j_{!*}\mathcal L[2]$, and
\begin{equation}\label{cycoll:eq:middle-stalks}
 i^*\mathcal I=V_1[1],\qquad i^!\mathcal I=V_2[-1].
\end{equation}
It also has the following proper geometric realization.
Let $\pi:\widetilde X=\operatorname{Bl}_0X\to X$ and
$q:\widetilde X\to\mathbb P^1$ be the line-bundle projection.
The local system descends to $\ell$ on
$Y=\mathbb P^1\setminus\{0,-1,\infty\}$.
For $j_Y:Y\hookrightarrow\mathbb P^1$, set
$\bar\ell=j_{Y,!}\ell=Rj_{Y,*}\ell$. Then
\begin{equation}\label{cycoll:eq:blowup}
 R\pi_*q^*\bar\ell[2]\simeq\mathcal I,
\end{equation}
with the isomorphism restricting to the identity on $U$.
\end{theorem}
\begin{proof}
Ordinary restriction of \eqref{cycoll:eq:middle} is
$\tau_{\leq1}K=V_1[-1]$. Extraordinary restriction is
$V_2[-3]$, because $i^!Rj_*\mathcal L=0$.
The shift by $2$ gives \eqref{cycoll:eq:middle-stalks}.
These restrictions satisfy the perverse bounds at the point,
and both have zero degree-zero cohomology. The punctured
branches are clean. Thus $\mathcal I$ is perverse, with no
subobject or quotient supported on the complement of $U$.
This proves the intermediate-extension assertion.

The trivial radial monodromy gives the descent to $Y$;
the three nontrivial puncture monodromies make $\bar\ell$
clean. Proper base change identifies the stalk of the left
side of \eqref{cycoll:eq:blowup} with
$R\Gamma(Y,\ell)[2]=H^1(Y,\ell)[1]=V_1[1]$.
Proper Verdier duality computes its costalk from the same
calculation for $\ell^\vee$. Poincare duality on the
punctured direction curve identifies
$H^1(Y,\ell^\vee)^\vee$ with $H_c^1(Y,\ell)$.
Clean extension identifies compactly supported and ordinary
cohomology here. Tensoring with the positive radial circle
class identifies $H^1(Y,\ell)$ with $V_2$. The costalk
is therefore $V_2[-1]$. Away from $0$, $\pi$ is an
isomorphism. The same intermediate-extension characterization
proves \eqref{cycoll:eq:blowup} and its uniqueness over $U$.
\end{proof}

The blowup separates the branches. Its exceptional direction
curve nevertheless retains their joint cohomology. This is
the geometric reason that replacing the three-line arrangement
by independent axes would change the coefficient.

\section{The critical direct image and its point constituent}
\label{cycoll:sec:critical}

The cubic cover also defines a critical coefficient. Put
\[
 B=\mathbb A^3_{u,v,w},\qquad
 W=\lambda\bigl(w^3-uv(u+v)\bigr),\qquad
 \pi_B:B\times\mathbb A^1_\lambda\longrightarrow B,
\]
and let $a:Z\hookrightarrow B$ be the zero locus of the
factor multiplying $\lambda$. Write $\phi_W$ for ordinary
vanishing cycles, and $\varphi_W=\phi_W[-1]$ for their
perverse normalization. All coefficients in the next formula
are rational, with Tate data forgotten.

\begin{theorem}\label{cycoll:thm:critical}
There is a deck-equivariant isomorphism
\begin{equation}\label{cycoll:eq:dimensional-reduction}
 R\pi_{B,!}\varphi_W\mathbb Q_{B\times\mathbb A^1}[4]
       \simeq a_*\mathbb Q_Z[2].
\end{equation}
After finite pushforward along $p:Z\to X$, scalar extension to
$\mathbb Q(\zeta)$, and projection to
a nontrivial deck character, the coefficient is
$\mathcal P=j_!\mathcal L[2]$. It is perverse and fits into
the nonsplit exact sequence
\begin{equation}\label{cycoll:eq:point-extension}
 0\longrightarrow i_*V_1\longrightarrow\mathcal P
       \longrightarrow\mathcal I\longrightarrow0,
 \qquad \dim V_1=1.
\end{equation}
Thus this critical direct image and the intermediate extension
have the same restriction to $U$ and different extensions
across its complement.
\end{theorem}
\begin{proof}
The potential is linear in $\lambda$. Dimensional reduction
for one linear fiber coordinate gives
\[
 R\pi_{B,!}\phi_W\mathbb Q[-1]\simeq a_*\mathbb Q_Z[-2]
\]
\cite[Theorem~A.1]{DavisonCriticalCoHA}.
Shifting by the ambient complex dimension $4$ gives
\eqref{cycoll:eq:dimensional-reduction}. The deck action
preserves $W$, so naturality of the comparison makes it
deck-equivariant. In particular its character summands
are preserved.

Over a point of $D$, the fiber of $p$ is the single
topological point $w=0$, fixed by the deck group. Finite
proper base change therefore gives zero ordinary boundary
stalks for any nontrivial character summand of
$Rp_*\mathbb Q_Z$. On $U$ that summand is $\mathcal L$.
The extension-by-zero adjunction is an isomorphism on every
stalk, identifying the summand with $j_!\mathcal L$.

The point stalk of $\mathcal P$ is zero. By
\eqref{cycoll:eq:clean-triangle}, its costalk is $K[1]$,
with cohomology in degrees $0,1$. The punctured branches
are clean. These restrictions prove perversity.
The two truncation triangles of $K$ give a triangle
\[
 j_!\mathcal L\longrightarrow F
       \longrightarrow i_*V_1[-1]\longrightarrow.
\]
After shifting by $2$ and rotating, perverse cohomology
gives \eqref{cycoll:eq:point-extension}.
If the sequence split, its middle term would have point
stalk $V_1\oplus V_1[1]$, whereas that stalk is zero.
The sequence is therefore nonsplit.
\end{proof}

The point constituent can be seen directly on $Z$. Let
$i_Z:\{0\}\hookrightarrow Z$ and
$j_Z:Z\setminus\{0\}\hookrightarrow Z$.
The surface is the affine cone over the projective cubic
\[
 E=\{w^3=uv(u+v)\}\subset\mathbb P^2.
\]
The cubic $E$ is smooth: the partial derivatives of
$w^3-uv(u+v)$ have no common projective zero.
Let $\operatorname{IC}_Z=j_{Z,!*}\mathbb Q[2]$ denote
the intersection-cohomology complex.

\begin{proposition}\label{cycoll:prop:cone}
There is an exact sequence of perverse sheaves
\begin{equation}\label{cycoll:eq:cone-extension}
 0\longrightarrow i_{Z,*}H^1(E,\mathbb Q)
       \longrightarrow\mathbb Q_Z[2]
       \longrightarrow\operatorname{IC}_Z\longrightarrow0.
\end{equation}
Each nontrivial deck character of $H^1(E)$ has dimension
one after adjoining $\zeta$. Finite pushforward and character
projection take \eqref{cycoll:eq:cone-extension} to
\eqref{cycoll:eq:point-extension}.
\end{proposition}
\begin{proof}
The link of the cone is the unit circle bundle of
$\mathcal O_E(-1)$, whose Euler number is $-3$.
Its Gysin sequence gives
$H^1(\operatorname{Link}(Z),\mathbb Q)=H^1(E,\mathbb Q)$.
The cone localization triangle computes the costalk of
$\mathbb Q_Z$ as the fiber of
$\mathbb Q\to R\Gamma(\operatorname{Link}(Z),\mathbb Q)$.
After shift by $2$ this costalk has no negative cohomology,
and its degree-zero cohomology is $H^1(E,\mathbb Q)$.
The ordinary stalk is $\mathbb Q[2]$.

These are the perverse bounds at the vertex; away from
the vertex the surface is smooth. There is no point
quotient, since the ordinary stalk has no degree-zero
cohomology. The maximal point subobject is the degree-zero
costalk, namely $i_{Z,*}H^1(E,\mathbb Q)$.
Taking its quotient removes the point subobject without
introducing a point quotient, so the quotient is
$\operatorname{IC}_Z$.

The deck quotient of $E$ is $\mathbb P^1$. Averaging over
the finite group identifies invariant cohomology with
the cohomology of the quotient, hence $H^1(E)$ has no
invariant part. It has rational dimension two, and the
two conjugate nontrivial characters each occur once.
Finite pushforward is perverse exact and commutes with
intermediate extension. This proves the last assertion.
\end{proof}

The inclusion $i_*V_1\to\mathcal P$ in
\eqref{cycoll:eq:point-extension} is a nonzero supported
morphism whose map on the ordinary point stalk is zero.
Indeed adjunction identifies it with the inclusion of
$H^0(i^!\mathcal P)=V_1$ into $i^!\mathcal P=K[1]$.
It is therefore detected by the costalk. This is why ordinary
stalk maps alone cannot determine a critical operation
at the collision.

\section{A coefficient concentrated at the deepest collision}
\label{cycoll:sec:deep}

The cubic construction extends to any number of insertions.
For $n\geq3$, remove common translation by setting $z_1=0$,
and define
\begin{equation}\label{cycoll:eq:vandermonde}
 \begin{gathered}
 X_n=\mathbb A^{n-1},\qquad N=\binom n2,\qquad
 Q_n=\prod_{1\leq i<j\leq n}(z_j-z_i),\\
 U_n=X_n\setminus V(Q_n),\qquad
 j:U_n\hookrightarrow X_n,\qquad i:\{0\}\hookrightarrow X_n.
 \end{gathered}
\end{equation}
The origin is the deepest collision, where all positions meet.
For $q\in k^\times$, let $L_n(q)$ be the pullback by
$Q_n:U_n\to\mathbb C^*$ of the rank-one local system with
positive meridian monodromy $q$. Every pair-collision
hyperplane consequently has the same monodromy $q$.

When $q^N=1$, this coefficient is a character summand of
the finite cover
\begin{equation}\label{cycoll:eq:cover}
 p_n:Z_n=\{w^N=Q_n\}\longrightarrow X_n.
\end{equation}
The coefficient field is assumed to contain the required
roots of unity. The character is faithful when $q$ has
order $N$, but faithfulness will not be necessary.

\begin{lemma}\label{cycoll:lem:proper-collisions}
Suppose
\begin{equation}\label{cycoll:eq:proper-condition}
 q^{\binom r2}\ne1\qquad(2\leq r<n).
\end{equation}
Then the normal coefficient link vanishes at every
nontrivial proper collision stratum. Hence
$j_!L_n(q)\to Rj_*L_n(q)$ is an isomorphism on
$X_n\setminus\{0\}$.
\end{lemma}
\begin{proof}
A collision stratum is specified by a partition into clusters
of equal positions. At a point of that stratum, differences
between distinct clusters are nonzero. Its normal arrangement
is the product of the arrangements internal to the clusters.
A cluster of size $r$ has an internal Vandermonde of degree
$\binom r2$. Its positive radial circle therefore acts by
$q^{\binom r2}$.

The radial circle complex contracts by the inverse of
$q^{\binom r2}-1$. Tensoring its contraction with the
remaining normal factors contracts the entire normal link.
Every nontrivial proper stratum has a cluster of size
$2\leq r<n$, so \eqref{cycoll:eq:proper-condition} applies.
The cone of $j_!\to Rj_*$ vanishes on these strata, proving
the assertion.
\end{proof}

For a primitive $N$th root, condition
\eqref{cycoll:eq:proper-condition} holds because
$0<\binom r2<N$. In particular the calculation includes
nonconsecutive clusters and simultaneous disjoint collisions.
At four insertions, for example, it includes $13\mid24$.
Only the coefficient link vanishes; the stratum itself
remains part of the configuration geometry.

\subsection{The direction space}

On $U_n$, set $t=z_2$ and $x_i=z_i/t$ for $i\geq3$.
Then
\begin{equation}\label{cycoll:eq:directions}
 U_n\simeq\mathbb C^*\times Y_n,\qquad
 Y_n=\{x_i\ne0,1,\ x_i\ne x_j\}.
\end{equation}
The marked points $0,1,x_3,\ldots,x_n,\infty$ identify
$Y_n$ with the moduli space $\mathcal M_{0,n+1}$ of ordered
distinct points on a projective line, modulo projective
automorphisms. Its complex dimension is $m=n-2$.
Since $Q_n=t^Nq_n(x)$, the radial monodromy is $q^N$.
For $q^N=1$, the coefficient in
\eqref{cycoll:eq:directions} is the external product of
the constant radial coefficient and a rank-one system
$\ell_n(q)$ on $Y_n$.

\begin{theorem}\label{cycoll:thm:deep-link}
Assume $q^N=1$ and \eqref{cycoll:eq:proper-condition}.
The direction cohomology is concentrated in degree $m=n-2$:
\begin{equation}\label{cycoll:eq:direction-rank}
 H^a(Y_n,\ell_n(q))=0\quad(a\ne m),\qquad
 \dim_k H^m(Y_n,\ell_n(q))=m!.
\end{equation}
The deepest link $K_n=i^*Rj_*L_n(q)$ has precisely
\begin{equation}\label{cycoll:eq:deep-link}
 H^{n-2}(K_n)=V_n,\qquad H^{n-1}(K_n)=W_n,
 \qquad \dim V_n=\dim W_n=(n-2)!.
\end{equation}
\end{theorem}
\begin{proof}
The smooth compactification
$\overline{\mathcal M}_{0,n+1}$ has normal-crossing boundary,
whose divisors are indexed by stable bipartitions of the
marked points \cite[\S2.6]{cycoll:Brown}.
Choose the side not containing $\infty$. It is a cluster
of $r$ finite marks, with $2\leq r\leq n-1$.
The Vandermonde winding around its collision divisor is
$\binom r2$. Renormalizing the coordinates to $z_1=0$,
$z_2=1$ changes this winding by a multiple of $N$, since
simultaneous rescaling multiplies $Q_n$ by an $N$th power.
Thus the boundary monodromy is $q^{\binom r2}$.

At an intersection of boundary divisors the local normal
link is a product of circles, with these nontrivial
monodromies. Its coefficient complex contracts. Extension
by zero to the compactification equals right direct image,
and properness gives
\begin{equation}\label{cycoll:eq:compact-comparison}
 H_c^a(Y_n,\ell_n(q))\xrightarrow{\sim}
 H^a(Y_n,\ell_n(q)).
\end{equation}

Successively forgetting a moving point expresses $Y_n$
as a tower of locally trivial fibrations by punctured
projective lines. Local trivializations follow by moving
the punctures in disjoint disks. Each fiber has the
homotopy type of a finite graph. The local-coefficient
spectral sequence, applied inductively also to the
finite-rank systems given by fiber cohomology, proves
vanishing above complex dimension $m$.
The same bound applies to the dual coefficient.
Poincare duality on the oriented real $2m$-manifold then
gives compact-support vanishing below $m$.
Together with \eqref{cycoll:eq:compact-comparison}, these
bounds prove concentration in degree $m$.

The fiber Euler characteristics in the tower are
$-1,-2,\ldots,-m$, so
$\chi(Y_n)=(-1)^m m!$. Euler characteristic with a rank-one
local system is the same as for the constant system, as
seen from a finite cellular complex. Concentration proves
the dimension in \eqref{cycoll:eq:direction-rank}.
Finally \eqref{cycoll:eq:directions} tensors this group
with $H^*(\mathbb C^*,k)$ in degrees $0,1$.
Positive real scaling identifies the complement with
its local link at the origin, proving
\eqref{cycoll:eq:deep-link}.
\end{proof}

The first ranks are $1,2,6,24$ for $n=3,4,5,6$.
The cubic cover is the case $n=3$.

\subsection{The critical extension in every arity}

Put $d=n-1$, $B_n=X_n\times\mathbb A^1_w$, and define
\[
 \mathcal W_n=\lambda(w^N-Q_n),\qquad
 \pi:B_n\times\mathbb A^1_\lambda\longrightarrow B_n,
 \qquad a:Z_n\hookrightarrow B_n.
\]
We continue to use $\varphi=\phi[-1]$ for perverse
vanishing cycles and omit Tate data.

\begin{theorem}\label{cycoll:thm:deep-critical}
Under the hypotheses of Theorem~\ref{cycoll:thm:deep-link},
dimensional reduction is deck-equivariant and gives
\begin{equation}\label{cycoll:eq:deep-dimred}
 R\pi_!\varphi_{\mathcal W_n}
       k_{B_n\times\mathbb A^1}[n+1]
       \simeq a_*k_{Z_n}[n-1].
\end{equation}
After pushforward by $p_n$ and projection to the character
with monodromy $q$, the coefficient is the perverse sheaf
$\mathcal P_n=j_!L_n(q)[d]$. It has a nonsplit exact sequence
\begin{equation}\label{cycoll:eq:deep-extension}
 0\longrightarrow i_*V_n\longrightarrow\mathcal P_n
       \longrightarrow j_{!*}L_n(q)[d]\longrightarrow0,
 \qquad \dim V_n=(n-2)!.
\end{equation}
\end{theorem}
\begin{proof}
The one-coordinate dimensional-reduction identity is
$R\pi_!\phi_{\mathcal W_n}k[-1]\simeq a_*k_{Z_n}[-2]$.
Shifting by the ambient dimension $n+1$ gives
\eqref{cycoll:eq:deep-dimred}; naturality gives deck
equivariance. Every fiber of $p_n$ over $Q_n=0$ is the
single topological point $w=0$. Since the selected
character is nontrivial, its ordinary boundary stalks
vanish. Its direct image is therefore $j_!L_n(q)$.

By Lemma~\ref{cycoll:lem:proper-collisions}, the difference
between the two open extensions is supported only at $0$:
\[
 j_!L_n(q)\longrightarrow Rj_*L_n(q)
       \longrightarrow i_*K_n\longrightarrow.
\]
Define
\[
 F_n=\operatorname{fib}\left(Rj_*L_n(q)
           \longrightarrow i_*W_n[-d]\right),\qquad
 \mathcal I_n=F_n[d].
\]
The arrow is the top-cohomology quotient on the point link.
The two restrictions are
$i^*\mathcal I_n=V_n[1]$ and
$i^!\mathcal I_n=W_n[-1]$.
Thus $\mathcal I_n$ is the intermediate extension.
The point stalk of $\mathcal P_n$ is zero and its costalk
is $K_n[d-1]$, whose cohomology is in degrees $0,1$.
The extension is clean at every other boundary stratum,
so $\mathcal P_n$ is perverse.

The truncation triangles of $K_n$ now give
$i_*V_n\to\mathcal P_n\to\mathcal I_n\to i_*V_n[1]$,
which is \eqref{cycoll:eq:deep-extension} in the perverse
heart. A splitting would give the nonzero point stalk
$V_n\oplus V_n[1]$ to $\mathcal P_n$. This contradicts
its zero stalk.
\end{proof}

\subsection{Keeping the monodromy fixed}

The preceding construction can choose a primitive
$\binom n2$th root for each $n$. That choice changes
the coefficient with the number of insertions.
For a fixed nontrivial $q$ of finite order $\ell$, set
\begin{equation}\label{cycoll:eq:first-arity}
 n(\ell)=\min\left\{r\geq2:\ell\mid\binom r2\right\}.
\end{equation}
The minimum exists, since $\ell\mid\binom{2\ell}{2}$.

\begin{corollary}\label{cycoll:cor:first-arity}
Among $n\geq3$, the coefficient $L_n(q)$ has vanishing
proper-collision links and a nonzero deepest link exactly
at $n=n(\ell)$. For $\ell=2,3,4,5,6,7,8$, these arities
are respectively $4,3,8,5,4,7,16$.
\end{corollary}
\begin{proof}
Before the first such arity the total radial monodromy
$q^{\binom n2}$ differs from $1$, so the deepest link
contracts. At the first arity, all proper cluster
monodromies are nontrivial and
Theorem~\ref{cycoll:thm:deep-link} applies.
At any larger arity choose a cluster of size $n(\ell)$
and keep the other points distinct. Its normal link is
the nonzero link just computed, so proper-collision
vanishing fails. Testing the intervening triangular
numbers gives the listed values.
\end{proof}

For example $q=-1$ gives the deepest-only phenomenon
at four insertions, although it is not a faithful
character of the sixfold cover. No fixed nontrivial
finite-order scalar character gives that phenomenon
in every arity. Compatibility between different arities
requires actual coefficient maps along substitutions.
The family of finite covers alone does not supply those maps.

\section{The geometric comparison still required}
\label{cycoll:sec:comparison}

The moment chart and the finite covers answer different
coefficient questions. The chart computes relation
directions in a derived flag fiber. The covers determine
constructible extensions and their point constituents.
Neither calculation identifies its coefficient with the
oriented direct image of a Hall correspondence.

\begin{problem}\label{cycoll:prob:hall-comparison}
Construct the oriented critical Hall or Hecke coefficient
map on the flag chart \eqref{cycoll:eq:moment-complex},
with its support maps and cohomological shift.
On the separated chamber, compare its coefficient
operation with $\mathcal L_s$ in \eqref{cycoll:eq:lift},
including the homotopy \eqref{cycoll:eq:weighted-homotopy}.
At the collision retain both odd directions through
the full supported complex
\eqref{cycoll:eq:supported-complex} before applying
the geometric direct image.

For three and more positions, determine the extension
of that same geometric coefficient across every collision
stratum. Compare its ordinary stalks, costalks and
attaching maps with the extensions calculated here.
In particular determine whether a critical point
constituent occurs and how the actual operation acts
on it. The comparison must include the nonconsecutive
collision $13\mid2$ and its common intersection with
the other pair collisions.
\end{problem}

A framing and its moving position connection must be
part of this map when framed coefficients are used.
The charge-zero scalar unit, the local vacuum under
the diagonal tensor and a functional defining a state
remain distinct data. Their behavior cannot be inferred
from a dimension or an Euler ratio. The examples in
this chapter specify which classes and supported maps
the geometric comparison must account for.
